% !TEX root = MM2025.tex


%%%%%%%%%%%%%%%%%%%%%%%%%%%%%%%%%%%%%%%%%%%%%%%%%%
\section{{\sectionnine}\label{sec:Conclusion}}
%%%%%%%%%%%%%%%%%%%%%%%%%%%%%%%%%%%%%%%%%%%%%%%%%%

In this paper, we studied the micro sources and macro consequences of the gender pay gap. To interpret the empirical fact that men are sorted into employers with higher pay, we developed an equilibrium model of firm pay, amenities, and employment. A notable feature of our model was that it allowed for the allocation of and transfers to workers of both genders to reflect compensating differentials \citep{Rosen1986}, taste-based discrimination \citep{becker1971}, and labor market frictions \citep{Manning2003}. We provided a constructive proof of identification of all model parameters based on linked employer-employee data. We then used the estimated framework to shed light on the structure of employer compensation by gender and to simulate counterfactual experiments, including equal-treatment policies.

It should be evident that our methodology and quantitative results have many interesting implications beyond the application to gender. Here, we focus on three. %
%
First, our empirical estimates suggest that men and women do not share a common ranking of employers. This suggests that not all workers climb the same job ladder. By allowing for separate job ladders by gender, we have taken a first step in this direction. Our methodology can be applied to study employer heterogeneity across other population groups, for example, by parental status or occupation.

Second, our finding of a significant role for compensating differentials suggests that measured inequality in labor market outcomes may overstate the true level of inequality. We documented stark differences in the structure of total compensation across men and women, which suggests it may be fruitful to revisit other distributional phenomena, such as cross-country differences in welfare and the uneven evolution of welfare within countries over time.

Third, employers in our framework have more than one margin to adapt to the economic environment. When considering the impact of equal-treatment policies, minimum wages, or income taxation, for example, amenities may move in the opposite direction from pay. Thus, the consequences of these policies for equity and efficiency may differ from what data on pay alone might suggest.

A cautionary note regarding the interpretation of our results is in order. Workers' preferences over amenities, such as those related to work-life balance, plausibly capture the induced demand resulting from household arrangements and societal norms that differ between men and women \citep{Goldin2014, Goldin2024, CubasJuhnSilos2021, CubasJuhnSilos2023}. In this sense, the gender differences we document may reflect not preferences but constraints imposed by nonmarket factors.
